\subsection{Derivation of the Monte Carlo Implementation of the Parametric G-Formula}
\label{app:gformula_derivation}

Here provide a brief derivation showing why forward simulation from
the disease-evolution model generates sample paths from the post-intervention
covariate distribution $f^g(\bar{L}_K)$, and why Monte Carlo averaging over such
trajectories yields a valid numerical implementation of the parametric g-formula.

\subsubsection{Factorization of the Post-Intervention Covariate Distribution}

Consider a deterministic dynamic treatment policy
$g = (g_0,\ldots,g_K)$, where at each time $t$ the treatment is assigned as
$A_t = g_t(\bar{L}_t,\bar{A}_{t-1})$.
Under the assumptions of consistency, sequential ignorability, and positivity,
the joint distribution of the covariate history under policy $g$ and no censoring,
$f^g(\bar{l}_K)$, is identified by the g-formula.

By repeated application of the chain rule and conditioning on survival and lack
of censoring up to each time point, this distribution admits the factorization
\begin{align}
f^g(\bar{l}_K)
&=
f(l_0)
\prod_{t=0}^{K-1}
f\!\left(
l_{t+1}
\,\middle|\,
\bar{l}_t,\,
\bar{a}^g_t,\,
\bar{C}_t = \bar{Y}_t = 0
\right),
\label{eq:fg_factorization}
\end{align}
where $\bar{a}^g_t = (a^g_0,\ldots,a^g_t)$ is the treatment history induced by policy
$g$. Sequential ignorability ensures that each conditional density in
\eqref{eq:fg_factorization} coincides with the corresponding conditional density
estimable from the observed data-generating process.

\subsubsection{G-Formula for Discrete-Time Survival}

For a discrete-time survival outcome, the probability of surviving through time
$K+1$ under policy $g$ can be written as
\begin{align}
P(Y^g_{K+1} = 0)
&=
\sum_{\bar{l}_K}
\left[
\prod_{t=0}^{K}
P\!\left(
Y_{t+1} = 0,\,
C_{t+1} = 0
\,\middle|\,
\bar{L}_t = \bar{l}_t,\,
\bar{A}_t = \bar{a}^g_t,\,
\bar{C}_t = \bar{Y}_t = 0
\right)
\right]
f^g(\bar{l}_K).
\label{eq:gformula_survival}
\end{align}

Substituting the factorization of $f^g(\bar{l}_K)$ from
\eqref{eq:fg_factorization} into \eqref{eq:gformula_survival} yields an expression
for $P(Y^g_{K+1} = 0)$ as an expectation with respect to the post-intervention
covariate distribution.

\subsubsection{Monte Carlo Approximation}

Equation \eqref{eq:gformula_survival} involves a summation over all possible
covariate trajectories $\bar{l}_K$, which is computationally infeasible in
all but trivial settings. However, because $f^g(\bar{l}_K)$ is fully specified
by the product of conditional transition densities in
\eqref{eq:fg_factorization}, we can generate independent draws
$\bar{l}^{(1)}_K,\ldots,\bar{l}^{(M)}_K \sim f^g(\bar{L}_K)$
by recursively simulating
\[
L_{t+1}^{(m)} \sim
f\!\left(
L_{t+1}
\,\middle|\,
\bar{L}^{(m)}_t,\,
\bar{A}^{g}_t,\,
\bar{C}_t = \bar{Y}_t = 0
\right),
\qquad
t = 0,\ldots,K-1,
\]
with $A_t = g_t(\bar{L}^{(m)}_t,\bar{A}^{(m)}_{t-1})$.

The g-formula integral can then be approximated by the Monte Carlo estimator
\begin{align}
P(Y^g_{K+1} = 0)
&\approx
\frac{1}{M}
\sum_{m=1}^{M}
\prod_{t=0}^{K}
P\!\left(
Y_{t+1} = 0,\,
C_{t+1} = 0
\,\middle|\,
\bar{L}_t = \bar{l}^{(m)}_t,\,
\bar{A}_t = \bar{a}^g_t,\,
\bar{C}_t = \bar{Y}_t = 0
\right),
\end{align}
which converges almost surely to the true g-formula estimand as
$M \to \infty$ by the law of large numbers.

Thus, forward simulation from the disease-evolution model under policy $g$,
followed by averaging of the corresponding survival probabilities, constitutes
a valid numerical implementation of the parametric g-formula.
